% SOURCE_STATEMENT: "Characterize when phase separation occurs in one-dimensional driven density-conserving systems and how this is related to traffic-jamming phase transitions, including a criterion in terms of the size dependence of steady-state currents."
% Preparation and final self-audit: CONTRIBUTING.md prompts C and D; source/audit artifacts retained in the local physics_statements_2026_10_11 batch.
% !TeX program = lualatex
\documentclass[11pt]{article}
\usepackage[margin=25mm]{geometry}
\usepackage{fontspec}
\setmainfont{Latin Modern Roman}
\usepackage{amsmath,amssymb,mathtools}
\usepackage{unicode-math}
\setmathfont{Latin Modern Math}
\usepackage{xurl,hyperref}
\hypersetup{colorlinks=true,urlcolor=blue}
\setcounter{secnumdepth}{0}
\setlength{\parindent}{0pt}
\setlength{\parskip}{5pt}
\setlength{\emergencystretch}{3em}
\newcommand{\sref}[2]{\hyperlink{source:#1:#2}{[S#2]}}
\newcommand{\cataloguescope}[1]{\paragraph{Recorded target.}#1}
\newcommand{\R}{\mathbb R}
\newcommand{\C}{\mathbb C}
\newcommand{\N}{\mathbb N}
\newcommand{\Z}{\mathbb Z}
\newcommand{\Q}{\mathbb Q}
\newcommand{\F}{\mathbb F}
\newcommand{\eps}{\varepsilon}
\begin{document}
\section{Current Criteria for Phase Separation in Driven Systems}\label{statement}

\subsection{Definitions and mathematical statement}

\paragraph{Shared notation.}$\mathbb N=\{1,2,\ldots\}$; the standard number fields have their usual meanings. All additional conventions are stated locally.


\paragraph{Definitions and conventions.} A one-dimensional driven conserving system has locally moving particles, a fixed total density, and a nonequilibrium stationary current. Phase separation means macroscopic domains with distinct densities rather than a homogeneous stationary state. If $J_\ell$ is the outgoing steady current of a domain of length $\ell$, its large-$\ell$ dependence can be compared with condensation criteria in an associated zero-range process. A jamming transition is the formation of a persistent macroscopic congested region in the traffic interpretation.

\paragraph{Source scope and qualifications.} The report describes a proposed correspondence and quantitative criterion in selected models. It does not define one class covering every driven system or prove necessity and sufficiency without assumptions; those choices cannot be invented from this programmatic extract.

\paragraph{Catalogue mathematical target.}
\paragraph{Catalogue wording retained for traceability.}
Characterize when phase separation occurs in one-dimensional driven density-conserving systems and how this is related to traffic-jamming phase transitions, including a criterion in terms of the size dependence of steady-state currents.

\paragraph{Saved literature context.} The supplied catalogue's assessment dated 2026-08-21 records: This is a contextual statement about one-dimensional phase separation and traffic jams rather than a precise theorem; the extracted item has no standalone resolved claim and is retained as open. The assessment is retained as dated provenance; it is not a new comprehensive literature or open-status certification.
\subsection{Short English statement}
Relate macroscopic phase separation and traffic jams to how stationary currents depend on domain length.
\subsection{Sources}
\begin{itemize}
\item[\hypertarget{source:statement:1}{S1}] ULAM.AI and the UnsolvedMath Contributors, UnsolvedMath, catalogue entry. Catalogue curation under CC BY 4.0; definitions and English exposition added locally.
\url{https://huggingface.co/datasets/ulamai/UnsolvedMath}.
\item[\hypertarget{source:statement:2}{S2}] Phase Transitions (2007), Oberwolfach Reports, DOI 10.4171/owr/2007/28, pp. 1575--1650; official report source and mathematical context.
\url{https://doi.org/10.4171/owr/2007/28}.
\item[\hypertarget{source:statement:3}{S3}] Official EMS Press PDF of the same Oberwolfach report; PDF pages 46--47, printed pages 1620--1621.
\url{https://ems.press/content/serial-article-files/46114}.
\item[\hypertarget{source:statement:4}{S4}] Ergodic theory of traffic jams; reference in the saved dated catalogue literature review.
\url{https://doi.org/10.1023/A:1012202706850}.
\end{itemize}
\label{end:statement}
\end{document}
